
This work established the first empirically validated classical variance baseline for neural network training stochasticity using N=30 independent random seeds on simple MLPs trained on MNIST and Fashion-MNIST. Variance measurements range from 0.04\% (MNIST) to 0.59\% (Fashion-MNIST, 2-layer), with all conditions showing statistically significant non-zero variance (p < 0.05).

The causal mechanism was validated through hypothesis decomposition. Different seeds produce independent initial weights (p < 0.000001), which lead to different optimization trajectories and final model states (final distances 22.7-27.3, loss CV 2-3\%). Test accuracy variance reflects the distribution of local minima quality.

Bootstrap stability analysis revealed that N=30 is sufficient for detecting non-zero variance but yields confidence interval widths of 93-110\%, exceeding the 50\% stability threshold. This finding establishes a detection-vs-precision boundary, refining existing sample size theory for neural network contexts.

Task-dependency analysis showed 9-$10\times$ variance scaling between easy (MNIST 0.04\%) and medium-difficulty (Fashion-MNIST 0.35-0.59\%) tasks, indicating that variance measurement is practical for tasks with 80-95\% baseline accuracy but ceiling-compressed above 95\%.

The work provides measurement infrastructure for uncertainty quantification research, enabling calibration of UQ methods against seed-based variance baselines. It establishes validated measurement protocols for reproducibility studies and identifies sample size requirements for variance detection (N=30) vs. precise estimation (N>50 recommended).

Limitations include scope restriction to simple MLPs, insufficient bootstrap precision at N=30, MNIST ceiling effects, and potential i.i.d. assumption violations. Future work includes N sensitivity analysis, statistical triangulation, task difficulty gradient experiments, and extension to complex architectures.
